\begin{corollary}[5.5.4]
\EGAReaderTarget{ega4.c04.cor.5-5-4}{IV.2.s5.5.4-fr}
Para todo anillo noetheriano $A$, se tiene
\begin{equation}
\dim(A[T_1,\ldots,T_r])=\dim(A)+r
\tag{5.5.4.1}\label{IV.2.s5.5.4.1-fr}
\end{equation}
($T_i$ indeterminadas).
\end{corollary}

Se observará, por el contrario, que existen ejemplos de anillos locales \emph{no
noetherianos} $A$ tales que $\dim(A)=1$ y $\dim(A[T])=3$ [30].

\begin{corollary}[5.5.5]
\label{IV.2.s5.5.5-fr}
Para todo preesquema localmente noetheriano $X$, la dimensión de
$X[T_1,\ldots,T_r]=X\otimes_{\mathbf Z}\mathbf Z[T_1,\ldots,T_r]$
($T_i$ indeterminadas) es $\dim(X)+r$.
\end{corollary}

Esto se sigue de (5.5.4) y de (\textbf 0, 14.1.7).

\begin{corollary}[5.5.6]
\label{IV.2.s5.5.6-fr}
Bajo las hipótesis de \textup{(5.5.3)}, sea $\mathfrak q$ un ideal primo de
$B=A[T]$ tal que $\mathfrak q\cap A=\mathfrak p$; si $k$ y $k'$ son los
cuerpos residuales de $A_{\mathfrak p}$ y $B_{\mathfrak q}$, respectivamente,
se tiene
\begin{equation}
\dim(A_{\mathfrak p})+1
=\dim(B_{\mathfrak q})+\operatorname{deg.tr}_k k'
\tag{5.5.6.1}\label{IV.2.s5.5.6.1-fr}
\end{equation}
\end{corollary}

\begin{proof}
\oldpage[IV-2]{96}

Si $\mathfrak q=\mathfrak pB$, se tiene, según (5.5.3.1),
$\dim(B_{\mathfrak q})=\dim(A_{\mathfrak p})$, y en este caso $k'=k(T)$,
por lo que la fórmula (5.5.6.1) queda verificada. En caso contrario,
$\dim(B_{\mathfrak q})=\dim(A_{\mathfrak p})+1$, y como $\mathfrak q$
corresponde a un ideal primo $\mathfrak q'\ne0$ de $k[T]$, $k'$ es una
extensión algebraica de $k$, por lo que se obtiene de nuevo la fórmula
(5.5.6.1).
\end{proof}

\begin{lemma}[5.5.7]
\label{IV.2.s5.5.7-fr}
Sean $A$ un anillo local íntegro noetheriano, $\mathfrak m$ su ideal maximal y
$k$ su cuerpo residual.
\begin{enumerate}
\item[\textup{(i)}] Para todo anillo íntegro $B$ que contenga a $A$, tal que $B=A[x]$
(para un $x\in B$), y todo ideal primo $\mathfrak q$ de $B$ tal que
$\mathfrak q\cap A=\mathfrak m$, se tiene
\begin{equation}
\dim(A)+\operatorname{deg.tr}_A B
\geq\dim(B_{\mathfrak q})+\operatorname{deg.tr}_k k'
\tag{5.5.7.1}\label{IV.2.s5.5.7.1-fr}
\end{equation}
denotando por $k'$ el cuerpo residual de $B_{\mathfrak q}$ y por
$\operatorname{deg.tr}_A B$ el grado de trascendencia del cuerpo de
fracciones de $B$ sobre el de $A$.
\item[\textup{(ii)}] Supongamos que, para todo ideal maximal $\mathfrak n$ de $A[T]$
tal que $\mathfrak n\cap A=\mathfrak m$, el anillo $(A[T])_{\mathfrak n}$
es catenario; entonces los dos miembros de \textup{(5.5.7.1)} son iguales.
\end{enumerate}
\end{lemma}

\begin{proof}
\begin{enumerate}
\item[(i)] Si $x$ es trascendente sobre el cuerpo de fracciones de $A$, se
tiene $\operatorname{deg.tr}_A B=1$ y los dos miembros de (5.5.7.1) son iguales
en virtud de (5.5.6). En caso contrario, se tiene $B=A[T]/\mathfrak p$, donde
$\mathfrak p$ es un ideal primo $\ne0$ de $A[T]$, tal que
$\mathfrak p\cap A=0$ puesto que $B$ contiene a $A$; por tanto se tiene
$\operatorname{ht}(\mathfrak p)=1$ en virtud de (5.5.3). El ideal
$\mathfrak q$ de $B$ es de la forma $\mathfrak n/\mathfrak p$, donde
$\mathfrak n\supset\mathfrak p$ es un ideal primo de $A[T]$ tal que
$\mathfrak n\cap A=\mathfrak m$, y se tiene
$B_{\mathfrak q}=(A[T])_{\mathfrak n}/\mathfrak p(A[T])_{\mathfrak n}$;
la fórmula (\textbf 0, 16.1.4.1), aplicada a
$X=\operatorname{Spec}((A[T])_{\mathfrak n})$ y a
$Y=\operatorname{Spec}(B_{\mathfrak q})$, da
\begin{equation}
\dim((A[T])_{\mathfrak n})
\geq\operatorname{ht}(\mathfrak p(A[T])_{\mathfrak n})
+\dim(B_{\mathfrak q})=1+\dim(B_{\mathfrak q})
\tag{5.5.7.2}\label{IV.2.s5.5.7.2-fr}
\end{equation}
porque $\operatorname{ht}(\mathfrak p(A[T])_{\mathfrak n})
=\operatorname{ht}(\mathfrak p)=1$ en virtud de la correspondencia biunívoca
entre los ideales primos de $A[T]$ contenidos en $\mathfrak n$ y los ideales
primos de $(A[T])_{\mathfrak n}$. Finalmente, la fórmula (5.5.6.1) da
\begin{equation}
\dim((A[T])_{\mathfrak n})
=\dim(A)+1\setmin\operatorname{deg.tr}_k k'
\tag{5.5.7.3}\label{IV.2.s5.5.7.3-fr}
\end{equation}
puesto que los cuerpos residuales de $B_{\mathfrak q}$ y de
$(A[T])_{\mathfrak n}$ son los mismos; por otra parte, se tiene entonces
$\operatorname{deg.tr}_A B=0$, lo que termina de probar (5.5.7.1).
\item[(ii)] Si $(A[T])_{\mathfrak n}$ es catenario, los dos miembros de
(5.5.7.2) son iguales (\textbf 0, 16.1.4), y por tanto también los dos miembros de
(5.5.7.1).
\end{enumerate}
\end{proof}

